% GENERATED FILE. Edit canonical lessons, then run npm run generate:books.
% canonical-source-hash: fnv1a64:06f7f18751f55641
% canonical-lessons: TE-C174-avatala, TE-C174-atu, TE-C174-itu, TE-C174-guha, TE-C174-loya, TE-R174-first-pass-words-from-chapters-166-169, TE-R174-second-pass-words-from-chapters-170-174

\chapter{On the other side, That way, This way, Cave, Valley}
\label{ch:te-on-the-other-side-that-way-this-way-cave-valley}
\hlchaptermodality{174}

\begin{chapteropening}
\textbf{By the end of this chapter:} \emph{I can say on the other side, that way, this way, a cave and a valley.}

Complete the last lesson of chapter 174: I can say on the other side, that way, this way, a cave and a valley.
\end{chapteropening}

\section[\texorpdfstring{\te{అవతల} (avatala)}{అవతల (avatala)}]{\te{అవతల} (avatala) — on the other side, beyond}

\label{lesson:TE-C174-avatala}



\begin{quote}
Before the new one: say the Telugu for beside, then the Telugu for a post office.
\end{quote}

\subsection*{You'll want to know: \te{అవతల}}
\textbf{\te{అవతల}} — \emph{avatala} — \textquotedblleft{}on the other side, beyond\textquotedblright{}.

\begin{grammarlens}[title={What you've built}]
One more word for this chapter.
\end{grammarlens}

\subsection*{Your turn}
\begin{itemize}
\raggedright
  \item \emph{Say it:} \emph{avatala}
  \item \emph{Say it:} \emph{avatala}, once more
  \item \emph{Say it:} say \emph{avatala}
  \item \emph{Recall:} say the Telugu for beside, then the Telugu for a post office, then say \emph{avatala} again
\end{itemize}

\begin{tcolorbox}[breakable,colback=teal!4,colframe=teal!35!black,title={Before you move on}]
What does \te{అవతల} mean? (\textquotedblleft{}On the other side, beyond\textquotedblright{}.) Say it once more.
\end{tcolorbox}
\section[\texorpdfstring{\te{అటు} (aṭu)}{అటు (aṭu)}]{\te{అటు} (aṭu) — that way, in that direction}

\label{lesson:TE-C174-atu}



\begin{quote}
Before the new one: say the Telugu for a post office, then the Telugu for on the other side.
\end{quote}

\subsection*{You'll want to know: \te{అటు}}
\textbf{\te{అటు}} — \emph{aṭu} — \textquotedblleft{}that way, in that direction\textquotedblright{}.

\begin{grammarlens}[title={What you've built}]
One more word for this chapter.
\end{grammarlens}

\subsection*{Your turn}
\begin{itemize}
\raggedright
  \item \emph{Say it:} \emph{aṭu}
  \item \emph{Say it:} \emph{aṭu}, once more
  \item \emph{Say it:} say \emph{aṭu}
  \item \emph{Recall:} say the Telugu for a post office, then the Telugu for on the other side, then say \emph{aṭu} again
\end{itemize}

\begin{tcolorbox}[breakable,colback=teal!4,colframe=teal!35!black,title={Before you move on}]
What does \te{అటు} mean? (\textquotedblleft{}That way, in that direction\textquotedblright{}.) Say it once more.
\end{tcolorbox}
\section[\texorpdfstring{\te{ఇటు} (iṭu)}{ఇటు (iṭu)}]{\te{ఇటు} (iṭu) — this way, in this direction}

\label{lesson:TE-C174-itu}



\begin{quote}
Before the new one: say the Telugu for on the other side, then the Telugu for that way.
\end{quote}

\subsection*{You'll want to know: \te{ఇటు}}
\textbf{\te{ఇటు}} — \emph{iṭu} — \textquotedblleft{}this way, in this direction\textquotedblright{}.

\begin{grammarlens}[title={What you've built}]
One more word for this chapter.
\end{grammarlens}

\subsection*{Your turn}
\begin{itemize}
\raggedright
  \item \emph{Say it:} \emph{iṭu}
  \item \emph{Say it:} \emph{iṭu}, once more
  \item \emph{Say it:} say \emph{iṭu}
  \item \emph{Recall:} say the Telugu for on the other side, then the Telugu for that way, then say \emph{iṭu} again
\end{itemize}

\begin{tcolorbox}[breakable,colback=teal!4,colframe=teal!35!black,title={Before you move on}]
What does \te{ఇటు} mean? (\textquotedblleft{}This way, in this direction\textquotedblright{}.) Say it once more.
\end{tcolorbox}
\section[\texorpdfstring{\te{గుహ} (guha)}{గుహ (guha)}]{\te{గుహ} (guha) — a cave}

\label{lesson:TE-C174-guha}



\begin{quote}
Before the new one: say the Telugu for that way, then the Telugu for this way.
\end{quote}

\subsection*{You'll want to know: \te{గుహ}}
\textbf{\te{గుహ}} — \emph{guha} — \textquotedblleft{}a cave\textquotedblright{}.

\begin{grammarlens}[title={What you've built}]
One more word for this chapter.
\end{grammarlens}

\subsection*{Your turn}
\begin{itemize}
\raggedright
  \item \emph{Say it:} \emph{guha}
  \item \emph{Say it:} \emph{guha}, once more
  \item \emph{Say it:} say \emph{guha}
  \item \emph{Recall:} say the Telugu for that way, then the Telugu for this way, then say \emph{guha} again
\end{itemize}

\begin{tcolorbox}[breakable,colback=teal!4,colframe=teal!35!black,title={Before you move on}]
What does \te{గుహ} mean? (\textquotedblleft{}A cave\textquotedblright{}.) Say it once more.
\end{tcolorbox}
\section[\texorpdfstring{\te{లోయ} (lōya)}{లోయ (lōya)}]{\te{లోయ} (lōya) — a valley}

\label{lesson:TE-C174-loya}



\begin{quote}
Before the new one: say the Telugu for this way, then the Telugu for a cave.
\end{quote}

\subsection*{You'll want to know: \te{లోయ}}
\textbf{\te{లోయ}} — \emph{lōya} — \textquotedblleft{}a valley\textquotedblright{}.

\begin{grammarlens}[title={What you've built}]
One more word for this chapter.
\end{grammarlens}

\subsection*{Your turn}
\begin{itemize}
\raggedright
  \item \emph{Say it:} \emph{lōya}
  \item \emph{Say it:} \emph{lōya}, once more
  \item \emph{Say it:} say \emph{lōya}
  \item \emph{Recall:} say the Telugu for this way, then the Telugu for a cave, then say \emph{lōya} again
\end{itemize}

\begin{tcolorbox}[breakable,colback=teal!4,colframe=teal!35!black,title={Before you move on}]
What does \te{లోయ} mean? (\textquotedblleft{}A valley\textquotedblright{}.) Say it once more.
\end{tcolorbox}
\section[(dialogue)]{Words from chapters 166-169 — a first pass}

\label{lesson:TE-R174-first-pass-words-from-chapters-166-169}



\begin{quote}
Say the words for a cave and a valley.
\end{quote}

\subsection*{Your turn}
Say these aloud:

\begin{itemize}
\raggedright
  \item to continue, to observe, to correct, to mutter, to separate
  \item a headache, diarrhoea, a tablet, sweat, a swelling
  \item jealousy, pride, enthusiasm, peace, patience
  \item a sore, a blister, a rash, a scar, a sprain
\end{itemize}

\begin{tcolorbox}[breakable,colback=teal!4,colframe=teal!35!black,title={Before you move on}]
To continue? (\textbf{\te{కొనసాగు}}, \emph{konasāgu}.) To observe? (\textbf{\te{పాటించు}}, \emph{pāṭiñcu}.) To correct? (\textbf{\te{దిద్దు}}, \emph{diddu}.) To mutter? (\textbf{\te{గొణుగు}}, \emph{goṇugu}.) To separate? (\textbf{\te{విడిపోవు}}, \emph{viḍipōvu}.) A headache? (\textbf{\te{తలనొప్పి}}, \emph{talanoppi}.) Diarrhoea? (\textbf{\te{విరేచనాలు}}, \emph{virēcanālu}.) A tablet? (\textbf{\te{మాత్ర}}, \emph{mātra}.) Sweat? (\textbf{\te{చెమట}}, \emph{cemaṭa}.) A swelling? (\textbf{\te{వాపు}}, \emph{vāpu}.) Jealousy? (\textbf{\te{అసూయ}}, \emph{asūya}.) Pride? (\textbf{\te{గర్వం}}, \emph{garvaṁ}.) Enthusiasm? (\textbf{\te{ఉత్సాహం}}, \emph{utsāhaṁ}.) Peace? (\textbf{\te{శాంతి}}, \emph{śānti}.) Patience? (\textbf{\te{సహనం}}, \emph{sahanaṁ}.) A sore? (\textbf{\te{పుండు}}, \emph{puṇḍu}.) A blister? (\textbf{\te{బొబ్బ}}, \emph{bobba}.) A rash? (\textbf{\te{దద్దుర్లు}}, \emph{daddurlu}.) A scar? (\textbf{\te{మచ్చ}}, \emph{macca}.) A sprain? (\textbf{\te{బెణుకు}}, \emph{beṇuku}.)
\end{tcolorbox}
\section[(dialogue)]{Words from chapters 170-174 — a second pass}

\label{lesson:TE-R174-second-pass-words-from-chapters-170-174}



\begin{quote}
Say the words for time and sometimes.
\end{quote}

\subsection*{Your turn}
Say these aloud:

\begin{itemize}
\raggedright
  \item time, sometimes, often, a date, a weekend
  \item sunset, every day, usually, rarely, constantly
  \item youth, old age, meanwhile, every year, forever
  \item a veranda, a backyard, a restaurant, beside, a post office
  \item on the other side, that way, this way, a cave, a valley
\end{itemize}

\begin{tcolorbox}[breakable,colback=teal!4,colframe=teal!35!black,title={Before you move on}]
Time? (\textbf{\te{సమయం}}, \emph{samayaṁ}.) Sometimes? (\textbf{\te{అప్పుడప్పుడు}}, \emph{appuḍappuḍu}.) Often? (\textbf{\te{తరచుగా}}, \emph{taracugā}.) A date? (\textbf{\te{తేదీ}}, \emph{tēdī}.) A weekend? (\textbf{\te{వారాంతం}}, \emph{vārāntaṁ}.) Sunset? (\textbf{\te{సూర్యాస్తమయం}}, \emph{sūryāstamayaṁ}.) Every day? (\textbf{\te{ప్రతిరోజు}}, \emph{pratirōju}.) Usually? (\textbf{\te{సాధారణంగా}}, \emph{sādhāraṇaṅgā}.) Rarely? (\textbf{\te{అరుదుగా}}, \emph{arudugā}.) Constantly? (\textbf{\te{నిరంతరం}}, \emph{nirantaraṁ}.) Youth? (\textbf{\te{యవ్వనం}}, \emph{yavvanaṁ}.) Old age? (\textbf{\te{వృద్ధాప్యం}}, \emph{vṛddhāpyaṁ}.) Meanwhile? (\textbf{\te{ఇంతలో}}, \emph{intalō}.) Every year? (\textbf{\te{ఏటా}}, \emph{ēṭā}.) Forever? (\textbf{\te{ఎప్పటికీ}}, \emph{eppaṭikī}.) A veranda? (\textbf{\te{వరండా}}, \emph{varaṇḍā}.) A backyard? (\textbf{\te{పెరడు}}, \emph{peraḍu}.) A restaurant? (\textbf{\te{హోటలు}}, \emph{hōṭalu}.) Beside? (\textbf{\te{పక్కన}}, \emph{pakkana}.) A post office? (\textbf{\te{పోస్టాఫీసు}}, \emph{pōsṭāphīsu}.) On the other side? (\textbf{\te{అవతల}}, \emph{avatala}.) That way? (\textbf{\te{అటు}}, \emph{aṭu}.) This way? (\textbf{\te{ఇటు}}, \emph{iṭu}.) A cave? (\textbf{\te{గుహ}}, \emph{guha}.) A valley? (\textbf{\te{లోయ}}, \emph{lōya}.)
\end{tcolorbox}
